\documentclass[Orbiter Developer Manual.tex]{subfiles}
\begin{document}

\section{Orbiter Source Code}

\subsection{Introduction}
As Orbiter is open-source, you can contribute to its development!
Below are instructions on how to download and setup and environment to build Orbiter by yourself.


\subsection{Setup}
Get the source repository of the Linux version of Orbiter from GitHub (\url{https://github.com/racerx2/orbiter-2027-linux}). The branch Custom holds Orbiter 2027 with the features of the Linux edition; main holds the line-by-line port of the Windows version

\begin{lstlisting}
git clone -b Custom git@github.com:racerx2/orbiter-2027-linux.git
\end{lstlisting}

\noindent
or

\begin{lstlisting}
git clone -b Custom https://github.com/racerx2/orbiter-2027-linux.git
\end{lstlisting}

\noindent
This is the Linux port of Orbiter; the Windows version is developed upstream at \url{https://github.com/orbitersim/orbiter}.\\
To compile Orbiter on Linux, you need CMake (3.28 or later), Ninja and GCC with C++20 support, and the development packages of Qt 6.5 or newer, Vulkan 1.3 (headers 1.3.246 or newer) with glslang, PipeWire, Wayland with wayland-protocols, libpng and zlib, as well as Python 3. The file COMPILE.md in the Orbiter sources lists the package names. The remaining libraries (Lua, Dear ImGui and others) are fetched by CMake.\\
The Linux version of Orbiter is a 64-bit (x86-64) application. To build it:

\begin{lstlisting}
cmake --preset linux-x64-release
cmake --build --preset linux-x64-release
ctest --preset linux-x64-release
\end{lstlisting}

\noindent
The binaries and data are placed in out/build/linux-x64-release, in the same layout as the installed Orbiter. Other presets are linux-x64-debug, linux-x64-asan (with address sanitizer) and linux-x64-tracy (with the Tracy profiler).\\
If you want to build the documentation (CMake option ORBITER\_MAKE\_DOC), you need a few additional tools:

\begin{itemize}
\item a LaTeX compiler suite such as TeX Live.
\item Doxygen and Graphviz for building the source-level documentation for developers.
\end{itemize}

\noindent
If the documentation build fails, try limiting the build to a single thread; some of the document compilers are not thread-safe.


\subsection{The release package}
The download of Orbiter 2027 for Linux is a portable package that runs on the common x86-64 distributions with glibc 2.34 or newer. The scripts that build it are in packaging/linux in the Orbiter sources:

\begin{itemize}
\item Dockerfile: the build environment, AlmaLinux 9 (glibc 2.34) with GCC 14 (gcc-toolset-14), CMake, Ninja and Qt 6.8 LTS. deps.sh builds Vulkan-Headers, glslang and Khronos' shader-object layer (VK\_LAYER\_KHRONOS\_shader\_object) into it.
\item build.sh <source> <out> <version> [manuals folder]: configures, builds and installs Orbiter from a copy of the source (without the high-resolution texture archives), copies the manuals, runs collect.py and packs Orbiter2027-Linux-x86\_64-<version>.tar.xz with its .sha256 file.
\item collect.py: copies Qt (with its plugins and the QML modules the Launchpad skins use) and every library outside the system set into lib/, sets each file's RUNPATH relative to lib/, writes qt.conf, and writes deps.list with only the system libraries, by package name for apt, dnf, zypper and pacman, which the OpenOrbiter launcher offers to install. It stops if a file needs a newer glibc (2.34) or libstdc++ (GLIBCXX 3.4.29, CXXABI 1.3.13) than the baseline.
\item test.sh <package> [images]: checks a package in containers of other distributions: the launcher's package names install, every library resolves, the launcher's checks pass, a scenario runs headless, the Launchpad loads a QML skin (offscreen and X11), and the VulkanClient renders on Mesa's software Vulkan driver (lavapipe), with and without the shader-object layer and under a headless Wayland compositor (weston).
\end{itemize}

\begin{lstlisting}
docker build -t orbiter-portable-build packaging/linux
docker run --rm -v $PWD:/src -v $PWD/out/portable:/out \
    orbiter-portable-build \
    /src/packaging/linux/build.sh /src /out <version>
packaging/linux/test.sh out/portable/Orbiter2027-Linux-x86_64-<version>.tar.xz
\end{lstlisting}

\noindent
The VulkanClient finds the shader-object layer in lib/vulkan of the package: it adds that folder to VK\_LAYER\_PATH (with the loader's own folders) before its first Vulkan call and enables the layer when it is there. The layer stays inactive on drivers that have VK\_EXT\_shader\_object themselves.


\subsection{Contribute}
You can contribute to the development of Orbiter, either by fixing a bug or by adding a new feature. You'll need to create your own fork of Orbiter in GitHub, put your changes in a branch, and then create a "Pull Request" (PR) so that you changes can be analysed by others, before they are merged into Orbiter.\\
If you find an issue in the Linux version of Orbiter, but you don't know how to fix it, you can open a ticket (\url{https://github.com/racerx2/orbiter-2027-linux/issues}) to allow others to fix it.


\end{document}
